\section{Quasi-coherent sheaves on Proj}
\label{section-quasi-coherent-proj}

\noindent
Let $S$ be a graded ring. Let $M$ be a graded $S$-module.
We saw in Lemma \ref{lemma-proj-sheaves} how to construct a quasi-coherent
sheaf of modules $\widetilde{M}$ on $\text{Proj}(S)$ and a map
\begin{equation}
\label{equation-map-global-sections}
M_0 \longrightarrow \Gamma(\text{Proj}(S), \widetilde{M})
\end{equation}
of the degree $0$ part of $M$ to the global sections of $\widetilde{M}$.
The degree $0$ part of the $n$th twist $M(n)$ of the graded module $M$ (see
Algebra, Section \ref{algebra-section-graded})
is equal to $M_n$. Hence we can get maps
\begin{equation}
\label{equation-map-global-sections-degree-n}
M_n \longrightarrow \Gamma(\text{Proj}(S), \widetilde{M(n)}).
\end{equation}
We would like to be able to perform this operation for any quasi-coherent
sheaf $\mathcal{F}$ on $\text{Proj}(S)$. We will do this by tensoring
with the $n$th twist of the structure sheaf, see
Definition \ref{definition-twist}. In order to relate the two notions
we will use the following lemma.

\begin{lemma}
\label{lemma-widetilde-tensor}
Let $S$ be a graded ring.
Let $(X, \mathcal{O}_X) = (\text{Proj}(S), \mathcal{O}_{\text{Proj}(S)})$
be the scheme of Lemma \ref{lemma-proj-scheme}.
Let $f \in S_{+}$ be homogeneous. Let $x \in X$ be a point
corresponding to the homogeneous prime $\mathfrak p \subset S$.
Let $M$, $N$ be graded $S$-modules.
There is a canonical map of $\mathcal{O}_{\text{Proj}(S)}$-modules
$$
\widetilde M \otimes_{\mathcal{O}_X} \widetilde N
\longrightarrow
\widetilde{M \otimes_S N}
$$
which induces the canonical map
$
M_{(f)} \otimes_{S_{(f)}} N_{(f)}
\to
(M \otimes_S N)_{(f)}
$
on sections over $D_{+}(f)$ and the canonical map
$
M_{(\mathfrak p)} \otimes_{S_{(\mathfrak p)}} N_{(\mathfrak p)}
\to
(M \otimes_S N)_{(\mathfrak p)}
$
on stalks at $x$. Moreover, the following diagram
$$
\xymatrix{
M_0 \otimes_{S_0} N_0 \ar[r] \ar[d] &
(M \otimes_S N)_0 \ar[d] \\
\Gamma(X, \widetilde M \otimes_{\mathcal{O}_X} \widetilde N) \ar[r] &
\Gamma(X, \widetilde{M \otimes_S N})
}
$$
is commutative where the vertical maps are given by
(\ref{equation-map-global-sections}).
\end{lemma}

\begin{proof}
To construct a morphism as displayed is the same as constructing
a $\mathcal{O}_X$-bilinear map
$$
\widetilde M \times \widetilde N
\longrightarrow
\widetilde{M \otimes_S N}
$$
see Modules, Section \ref{modules-section-tensor-product}.
It suffices to define this on sections over the opens $D_{+}(f)$
compatible with restriction mappings. On $D_{+}(f)$ we use the
$S_{(f)}$-bilinear map
$M_{(f)} \times N_{(f)} \to (M \otimes_S N)_{(f)}$,
$(x/f^n, y/f^m) \mapsto (x \otimes y)/f^{n + m}$. Details omitted.
\end{proof}

\begin{remark}
\label{remark-not-isomorphism}
In general the map constructed in Lemma \ref{lemma-widetilde-tensor}
above is not an isomorphism. Here is an example. Let $k$
be a field. Let $S = k[x, y, z]$ with $k$ in degree $0$ and
$\deg(x) = 1$, $\deg(y) = 2$, $\deg(z) = 3$.
Let $M = S(1)$ and $N = S(2)$, see
Algebra, Section \ref{algebra-section-graded}
for notation. Then $M \otimes_S N = S(3)$.
Note that
\begin{eqnarray*}
S_z
& = &
k[x, y, z, 1/z] \\
S_{(z)}
& = &
k[x^3/z, xy/z, y^3/z^2]
\cong
k[u, v, w]/(uw - v^3) \\
M_{(z)} & = & S_{(z)} \cdot x + S_{(z)} \cdot y^2/z \subset S_z \\
N_{(z)} & = & S_{(z)} \cdot y + S_{(z)} \cdot x^2 \subset S_z \\
S(3)_{(z)} & = & S_{(z)} \cdot z \subset S_z
\end{eqnarray*}
Consider the maximal ideal $\mathfrak m = (u, v, w) \subset S_{(z)}$.
It is not hard to see that both $M_{(z)}/\mathfrak mM_{(z)}$
and $N_{(z)}/\mathfrak mN_{(z)}$ have dimension $2$ over
$\kappa(\mathfrak m)$. But
$S(3)_{(z)}/\mathfrak mS(3)_{(z)}$ has dimension $1$.
Thus the map $M_{(z)} \otimes N_{(z)} \to S(3)_{(z)}$ is not
an isomorphism.
\end{remark}









